\section{Introduction}
\label{sec:introduction}

A finite-dimensional Markovian open quantum system evolves under a Gorini--Kossakowski--Sudarshan--Lindblad (GKSL) generator. Its exact propagator is completely positive and trace preserving (CPTP) at every time \cite{Lindblad1976,GoriniKossakowskiSudarshan1976,BreuerPetruccione2002}. Complete positivity is stronger than trace preservation or positivity on one selected trajectory: the map must remain positive after extension by an arbitrary ancilla, equivalently its Choi matrix must be positive semidefinite \cite{Choi1975CP,Kraus1971,Stinespring1955}. This distinction is operational whenever a numerical update is reused for different inputs or applied to one subsystem of a larger register.

Standard direct RK schemes are not CPTP-preserving in general, and explicit RK formulas supply no generator-independent complete-positivity guarantee \cite{RieschJirauschek2019,AppeloCheng2025}. Numerical approaches to Lindblad evolution have a longer history: high-order stochastic unravelings, deterministic and adaptive low-rank approximations, and locally purified tensor networks predate the current wave of CPTP-by-construction methods \cite{SteinbachGarrawayKnight1995,LeBrisRouchon2013,LeBrisRouchonRoussel2015,WernerEtAl2016}. Recent Kraus-form, quantum-channel, exponential, low-rank, nested-Picard, and tensor-compressed schemes instead alter the update so that positivity or complete positivity is built into the method \cite{RieschPiklJirauschek2020,CaoLu2025,AppeloCheng2025,RobinRouchonSellem2025,ChenEtAl2026,GodinezRamirez2025,HuAppeloCheng2025NPI,DelMastroAppeloCheng2026TT,HuAppeloCheng2026Gregory}. Complementary work addresses deterministic propagation, a posteriori error control, randomized simulation, and optimized unraveling \cite{GuZhang2025,TureHyeonJang2026,EtienneyRobinRouchon2026,ChenLiLuYing2025,CaoHeLi2026Unraveling}.

We ask the complementary question without changing the RK formula: for a fixed generator and a specified numerical candidate construction, which finite-step maps are physical? The answer is not determined by convergence order, scalar absolute stability, componentwise positivity, absolute monotonicity, or strong-stability preservation \cite{BolleyCrouzeix1978,Kraaijevanger1986Absolute,Kraaijevanger1991Contractivity,Horvath1998Positivity,Horvath2005Threshold,GottliebShuTadmor2001,FerracinaSpijker2004,Higueras2005,BlanesIserlesMacnamara2022}. Complete positivity couples distinct Liouvillian modes through a channel inequality, so every scalar mode can be stable while the induced map lies outside the CPTP set.

The exact laboratory is the autonomous phase-covariant qubit family: coherent $z$-precession, downward and upward relaxation, and pure dephasing. The underlying dynamics and non-unital qubit-channel geometry are established \cite{FujiwaraAlgoet1999,RuskaiSzarekWerner2002,BraunEtAl2014,Hall2008Complete,Lankinen2016Complete,Filippov2020Phase,Siudzinska2022Mixtures,Siudzinska2023Geometry,LonigroChruscinski2023}. Pulling that geometry back through an RK stability function produces a one-dimensional semialgebraic CPTP stability set along each generator ray. To our knowledge, prior work has not classified its connected, reentrant, isolated, and detached components for an unchanged RK formula, nor the consequences of constructing several different candidate maps over the same interval.

On the laboratory-frame one-way, zero-dephasing ray, classical RK4 provides a closed-form topology. In the laboratory-frame one-way problem its CPTP set changes repeatedly as $\varpi=|\omega|/\Gamma$ varies. A narrow reentrant band occurs for $0.82736\ldots<\varpi<\sqrt3/2$; above a second bifurcation the positive admissible interval is detached from the origin. At high frequency its dimensional edges satisfy
\begin{equation}
 h_-(\varpi)=\frac{3\Gamma}{\omega^2}+O\!\left(\frac{\Gamma^3}{\omega^4}\right),
 \qquad
 h_+(\varpi)=\frac{2\sqrt2}{|\omega|}+O\!\left(\frac{\Gamma}{\omega^2}\right).
 \label{eq:intro_window}
\end{equation}
This topology is a property of the discretization--representation pair. An exact rotating frame removes the coherent frequency from the commuting benchmark. In the displayed transverse-field family, separate exact calculations at $\Omega_x/\Gamma\in\{0,0.05,0.10,0.20\}$ each retain two map-level CP components; these four pointwise certificates do not establish uniform persistence on the intervening interval.

Adaptive workflows add a second, more general source of structure. Step doubling constructs a coarse full-step map and a fine composition of two half steps; local extrapolation may form a negative-weight affine combination \cite{Shampine1985Doubling,NussleinRanochaKetcheson2021,Szyszkowicz1985Richardson,ZlatevEtAl2017Richardson,HavasiKazemi2018Richardson}. These are distinct linear maps over one interval. At $\varpi=2$ the nontrivial positive-step CPTP components of the full and two-half-step candidates are disjoint, so no positive total step makes both candidates physical. The Richardson extrapolate built from RK4 is an order-five, eighth-degree stability function whose first exponential defect points out of the CPTP set. A global Sturm certificate then proves non-CPTP behavior throughout a positive interval on which both constituent candidates are CPTP. Unlike the frequency-induced laboratory-frame topology, this nonrotating extrapolation failure is unaffected by an exact rotating-frame change.

Direct RK remains relevant because it is widely available, works naturally with sparse or operator-form Liouvillian actions, and is embedded in general scientific ODE workflows \cite{PuzzuoliEtAl2023}. Time-dependent or noncommuting driven generators may not admit one static frame that removes all fast terms at every stage. Frame transformations, integrating factors, splitting, and CPTP-by-construction propagators provide alternative numerical representations. Exact knowledge of the unchanged method's topology informs diagnostics, frame selection, component projection, and fallback decisions.

The exact Choi-block criterion holds for the full autonomous phase-covariant family. The six-regime classification is restricted to the laboratory-frame one-way, zero-dephasing ray. We next test whether the disconnected topology persists under local buffers and a noncommuting transverse field. We then analyze the full-step, equal-substep, mixture, and Richardson candidates, obtaining the disjoint full-step/two-half-step result and the global non-CPTP interval for the RK4 Richardson candidate.

The controller uses three outcomes---\textsc{pass}, \textsc{fail}, and \textsc{uncertain}---and certifies the exact binary64 step together with the construction rule for the candidate actually executed. A rejection caused only by complete positivity leaves the proportional--integral (PI) controller's error memory unchanged. The controller projects in the buffered and near-saddle tests and falls back in the exact-boundary test.

Sections~\ref{sec:criterion}--\ref{sec:rk4topology} derive the criterion and boundary topology. The remaining sections treat robustness, candidate maps, certified step control, operational consequences, and structure-preserving alternatives.
